\subsection{Sprint 4 --- 27/10/2023 até 17/11/2023}

Após a retrospectiva da Sprint 3, iniciamos a quinta e última sprint do projeto. O progresso alcançado anteriormente aumentou a motivação da equipe e criou uma expectativa positiva para a entrega final, inclusive com o desejo de disputar o prêmio de melhor projeto.

Nesta sprint, além de tratar pequenas dívidas técnicas acumuladas nas etapas anteriores, incorporamos a última \textit{user story} solicitada pelos stakeholders. Essa demanda envolvia a implementação da exportação de dados em formato \texttt{.csv}, contendo as informações processadas pela aplicação. A sprint teve foco predominante em backend e correção de bugs, buscando garantir estabilidade e funcionamento completo antes da entrega.

Durante essa etapa, minha atuação foi mais voltada a suporte do que ao desenvolvimento direto de novas funcionalidades. Contribuí em alguns ajustes de frontend, especialmente relacionados à definição final do design, mas minha principal responsabilidade foi auxiliar colegas da AGES I na configuração do ambiente de trabalho e do banco de dados. Dediquei tempo a orientar outros membros para que pudessem trabalhar com mais autonomia e eficiência.

No encerramento da sprint, acompanhei colegas mais experientes lidando com questões relacionadas ao Firebase, o que me proporcionou uma base inicial para explorar tecnologias semelhantes no futuro. Mesmo sem assumir grandes tarefas novas de implementação, considero que essa atuação de apoio foi importante para a equipe e para meu desenvolvimento, pois reforçou a importância de compartilhar conhecimento e ajudar a destravar o trabalho coletivo.

De modo geral, esta sprint foi mais leve do que as anteriores. Isso ocorreu principalmente porque o projeto já estava bem encaminhado, reduzindo a pressão por entregas de última hora. Além disso, o final do semestre e a proximidade das férias contribuíram para um clima mais tranquilo dentro da equipe.

O projeto foi encerrado em 17/11/2023, durante a quinta e última reunião com os stakeholders. A reunião foi positiva, e o projeto foi aceito por completo. No entanto, houve uma situação inesperada: fomos informados de que Cassio e Letícia, dois dos stakeholders com quem havíamos trabalhado, haviam sido desligados da Triider. Ainda assim, ambos participaram da reunião, demonstrando respeito e apreço pelo trabalho desenvolvido pela equipe.

A Sprint 4 encerrou o projeto de forma positiva, embora com essa reviravolta inesperada. A equipe pôde concluir a entrega com mais tranquilidade, apoiada pelo avanço substancial das sprints anteriores. Para mim, essa fase reforçou que contribuir para um projeto não significa apenas implementar funcionalidades, mas também apoiar colegas, compartilhar conhecimento e colaborar para que o produto chegue ao fim com qualidade.
